% SPDX-FileCopyrightText: 2024 IObundle
%
% SPDX-License-Identifier: MIT

\documentclass{ug}
\usepackage{xltabular}

% SPDX-FileCopyrightText: 2024 IObundle
%
% SPDX-License-Identifier: MIT

\title{%
\Huge IOb-Cache \\
 \vspace*{3cm}
\Large A Configurable Cache
}

\header{IOb-Cache, a Configurable Cache}

\date{\today}
\category{User Guide, \input{\NAME_version.tex}, Build unknown}

% SPDX-FileCopyrightText: 2024 IObundle
%
% SPDX-License-Identifier: MIT

\definecolor{iob-green}{rgb}{0.0,1.0,0.80}
\definecolor{iob-blue}{rgb}{0.90196,0.94902,1}
\definecolor{dark-blue}{RGB}{155,194,230}
\definecolor{mid-blue}{RGB}{189,215,238}
\definecolor{light-blue}{RGB}{221,235,247}


\begin{document}

\maketitle
\pagenumbering{gobble}

\vspace*{\fill}
User Guide, \input{\NAME_version.tex}, Build unknown
\hspace*{\fill} \includegraphics[keepaspectratio,scale=.7]{Logo.png}

\cleardoublepage
\pagenumbering{roman}
\setcounter{page}{1}
% SPDX-FileCopyrightText: 2024 IObundle
%
% SPDX-License-Identifier: MIT

\subsection*{Document Version History}

\begin{table}[h]
  \centering
  \begin{tabularx}{\textwidth}{|c|c|c|X|}
    \hline
    \rowcolor{iob-green}
             {\bf Version} & {\bf Date} & {\bf Person} & {\bf Changes from previous version} \\ \hline
             \input ./revhist_desc
  \end{tabularx}
\end{table}

\cleardoublepage
\tableofcontents
\clearpage
\listoftables
\clearpage
\listoffigures
\cleardoublepage
\pagenumbering{arabic}
\setcounter{page}{1}
\section{Introduction}
\label{sec:intro}
% SPDX-FileCopyrightText: 2024 IObundle
%
% SPDX-License-Identifier: MIT

IOb-Cache is an open-source configurable pipelined memory cache. The
processor-side interface (front-end) uses IObundle's Native Pipelined Interface
(NPI). The memory-side interface (back-end) can also be configured to use NPI or
the widely used AXI4 interface. The address and data widths of the front-end and
back-end are configurable to support multiple user cores and memories. IOb-Cache
is a K-Way Set-Associative cache, where K can vary from 1 (directly mapped) to 8
or more ways, provided the operating frequency after synthesis is
acceptable. IOb-Cache supports the two most common write policies: Write-Through
Not-Allocate and Write-Back Allocate.

IOb-Cache was developed in the scope of João Roque's master's thesis in
Electrical and Computer Engineering at the Instituto Superior Técnico of the
University of Lisbon. The Verilog code works well in IObundle's IOb-SoC system
(https://github.com/IObundle/iob-soc) both in simulation and FPGA. To be used in
an ASIC, it would need to be lint-cleaned and verified more thoroughly by RTL
simulation to achieve 100\% code coverage desirably.

% SPDX-FileCopyrightText: 2024 IObundle
%
% SPDX-License-Identifier: MIT

\begin{figure}[!htbp]
  \centerline{\includegraphics[width=14cm]{symb.pdf}}
  \vspace{0cm}
  \caption{IP Core Symbol}
  \label{fig:symbol}
\end{figure}


\subsection{Features}
\label{sec:feat}
% SPDX-FileCopyrightText: 2024 IObundle
%
% SPDX-License-Identifier: MIT

\begin{itemize}
  \itemsep-0.5em
\item Pipelined operation allowing consecutive one-cycle reads and writes
\item IObundle's Native Pipelined Interface (NPI) front-end native interface on the processor side (front-end)
\item NPI or AXI4 interface on the memory side (back-end)
\item Configurable address and data widths on the front-end and back-end interfaces for supporting a variety of different systems
\item Configurable number of lines and words per line
\item Configurable K-Way Set-Associativity ($k \geq 1$)
\item Configurable line replacement policy: LRU, MRU-based PLRU, and tree-based PLRU.
\item Configurable Write-Through Not-Allocate and Write-Back Allocate policies
\item Configurable Write-Through buffer depth
\item Optional control address space for cache invalidation, accessing the write through buffer status and read/write hit/miss counters
\end{itemize}


\subsection{Benefits}
\label{sec:benef}
% SPDX-FileCopyrightText: 2024 IObundle
%
% SPDX-License-Identifier: MIT

\begin{itemize}
  \itemsep-0.5em
\item Easy to integrate hardware and bare-metal software
\item Low usage of FPGA resources or silicon area 
\item Operating frequency suitable for low-cost FPGAs and less recent ASIC nodes
\item Low power consumption
\end{itemize}


\subsection{Deliverables}
\label{sec:deliv}
% SPDX-FileCopyrightText: 2024 IObundle
%
% SPDX-License-Identifier: MIT

\begin{itemize}                                            
  \itemsep-0.5em
\item Verilog source code and testbench
\item Verilator testbench
\item FPGA implementation scripts
\item Bare-metal software driver                                                   
\item Documentation Latex sources
\item Example System on Chip using IOb-Cache
\end{itemize}   


\ifdefined\SECTIONCLEARPAGE
\clearpage
\fi
\section{Description}

This section gives a detailed description of the IP core. The high-level block
diagram is presented, along with a description of its blocks. The parameters and
macros that define the core configuration are listed and explained. The
interface signals are enumerated and described; if timing diagrams are needed,
they are shown after the interface signals. Finally the Control and Status
Registers (CSR) are outlined and explained.

\subsection{Block Diagram}
\label{sec:bdd}
% SPDX-FileCopyrightText: 2024 IObundle
%
% SPDX-License-Identifier: MIT

Figure~\ref{fig:bd} presents a high-level block diagram of the core, followed by
a brief description of each block.

% SPDX-FileCopyrightText: 2024 IObundle
%
% SPDX-License-Identifier: MIT

\begin{figure}[H]
  \centering {\includegraphics[width=\columnwidth]{bd.pdf}}
  \vspace{-0.7cm}
  \caption{High-Level Block Diagram}
  \label{fig:bd}
\end{figure}


\setlist[description]{font=\bfseries\MakeUppercase, labelindent=\parindent, leftmargin=*}

The Verilog modules in the top-level entity of the core are         described in the following tables. The table elements represent         the blocks in the Block Diagram.

\begin{table}[H]
  \centering
  \begin{tabularx}{\textwidth}{|l|l|X|}

    \hline
    \rowcolor{iob-green}
    {\bf Module} & {\bf Name} & {\bf Description}  \\ \hline \hline

    \input general_operation_blocks_tab

  \end{tabularx}
  \caption{General operation group}
  \label{general_operation_blocks_tab:is}
\end{table}
\clearpage


\ifdefined\CONFS
\subsection{Configuration}
\label{sec:ipconfig}

The following tables describe the IP core configuration. The core may be configured using macros or parameters:

\begin{description}
\item \textbf{'M'} Macro: a Verilog macro or ``\`define'' directive is used to include or exclude code segments, to create core configurations that are valid for all instances of the core.
\item \textbf{'P'} Parameter: a Verilog parameter is passed to each instance of the core and defines the configuration of that particular instance.
\end{description}

\begin{table}[H]
  \centering
  \begin{tabularx}{\textwidth}{|l|c|c|c|c|X|}

    \hline
    \rowcolor{iob-green}
    {\bf Configuration} & {\bf Type} & {\bf Min} & {\bf Typical} & {\bf Max} & {\bf Description} \\ \hline \hline

    \input general_operation_confs_tab

  \end{tabularx}
  \caption{General operation group}
  \label{general_operation_confs_tab:is}
\end{table}
\clearpage
\fi

\subsection{Interface Signals}
\label{sec:ifsig}
The interface signals of the core are described in the following tables.

\begin{table}[H]
  \centering
  \begin{tabularx}{\textwidth}{|l|l|r|X|}

    \hline
    \rowcolor{iob-green}
    {\bf Name} & {\bf Direction} & {\bf Width} & {\bf Description}  \\ \hline \hline

    \input clk_en_rst_s_if_tab

  \end{tabularx}
  \caption{Clock, clock enable and reset}
  \label{clk_en_rst_s_if_tab:is}
\end{table}

\begin{table}[H]
  \centering
  \begin{tabularx}{\textwidth}{|l|l|r|X|}

    \hline
    \rowcolor{iob-green}
    {\bf Name} & {\bf Direction} & {\bf Width} & {\bf Description}  \\ \hline \hline

    \input iob_s_if_tab

  \end{tabularx}
  \caption{Front-end interface}
  \label{iob_s_if_tab:is}
\end{table}

\begin{table}[H]
  \centering
  \begin{tabularx}{\textwidth}{|l|l|r|X|}

    \hline
    \rowcolor{iob-green}
    {\bf Name} & {\bf Direction} & {\bf Width} & {\bf Description}  \\ \hline \hline

    \input iob_m_if_tab

  \end{tabularx}
  \caption{Back-end interface}
  \label{iob_m_if_tab:is}
\end{table}

\begin{table}[H]
  \centering
  \begin{tabularx}{\textwidth}{|l|l|r|X|}

    \hline
    \rowcolor{iob-green}
    {\bf Name} & {\bf Direction} & {\bf Width} & {\bf Description}  \\ \hline \hline

    \input axi_m_if_tab

  \end{tabularx}
  \caption{AXI4 interface}
  \label{axi_m_if_tab:is}
\end{table}

\begin{table}[H]
  \centering
  \begin{tabularx}{\textwidth}{|l|l|r|X|}

    \hline
    \rowcolor{iob-green}
    {\bf Name} & {\bf Direction} & {\bf Width} & {\bf Description}  \\ \hline \hline

    \input axi_write_m_if_tab

  \end{tabularx}
  \caption{AXI4 write interface}
  \label{axi_write_m_if_tab:is}
\end{table}

\begin{table}[H]
  \centering
  \begin{tabularx}{\textwidth}{|l|l|r|X|}

    \hline
    \rowcolor{iob-green}
    {\bf Name} & {\bf Direction} & {\bf Width} & {\bf Description}  \\ \hline \hline

    \input axi_read_m_if_tab

  \end{tabularx}
  \caption{AXI4 read interface}
  \label{axi_read_m_if_tab:is}
\end{table}

\begin{table}[H]
  \centering
  \begin{tabularx}{\textwidth}{|l|l|r|X|}

    \hline
    \rowcolor{iob-green}
    {\bf Name} & {\bf Direction} & {\bf Width} & {\bf Description}  \\ \hline \hline

    \input fe_io_if_tab

  \end{tabularx}
  \caption{Front-end interface (IOb native slave)}
  \label{fe_io_if_tab:is}
\end{table}

\begin{table}[H]
  \centering
  \begin{tabularx}{\textwidth}{|l|l|r|X|}

    \hline
    \rowcolor{iob-green}
    {\bf Name} & {\bf Direction} & {\bf Width} & {\bf Description}  \\ \hline \hline

    \input be_io_if_tab

  \end{tabularx}
  \caption{Back-end interface}
  \label{be_io_if_tab:is}
\end{table}

\begin{table}[H]
  \centering
  \begin{tabularx}{\textwidth}{|l|l|r|X|}

    \hline
    \rowcolor{iob-green}
    {\bf Name} & {\bf Direction} & {\bf Width} & {\bf Description}  \\ \hline \hline

    \input ie_io_if_tab

  \end{tabularx}
  \caption{Cache invalidate and write-trough buffer IO chain}
  \label{ie_io_if_tab:is}
\end{table}

\begin{table}[H]
  \centering
  \begin{tabularx}{\textwidth}{|l|l|r|X|}

    \hline
    \rowcolor{iob-green}
    {\bf Name} & {\bf Direction} & {\bf Width} & {\bf Description}  \\ \hline \hline

    \input ge_i_if_tab

  \end{tabularx}
  \caption{General Interface Signals}
  \label{ge_i_if_tab:is}
\end{table}
\clearpage

%timing diagrams
\ifdefined\TD
\subsection{Timing Diagrams}
\label{sec:td}
\input{td}
\fi

%software components
\ifdefined\CSRS
\subsection{Control and Status Registers}
\label{sec:csrs}

    The software accessible registers of the core are described in the following
    tables. The tables give information on the name, read/write capability, address, width in bits, and a textual description.

    \begin{table}[H]
      \centering
      \begin{tabularx}{\textwidth}{|l|c|c|c|c|X|}
        
        \hline
        \rowcolor{iob-green}
        {\bf Name} & {\bf R/W} & {\bf Addr} & {\bf Width} & {\bf Default} & {\bf Description} \\ \hline

        \input cache_csrs_tab
     
      \end{tabularx}
      \caption{CACHE software accessible registers.}
      \label{cache_csrs_tab:is}
    \end{table}

    \begin{table}[H]
      \centering
      \begin{tabularx}{\textwidth}{|l|c|c|c|c|X|}
        
        \hline
        \rowcolor{iob-green}
        {\bf Name} & {\bf R/W} & {\bf Addr} & {\bf Width} & {\bf Default} & {\bf Description} \\ \hline

        \input general_operation_csrs_tab
     
      \end{tabularx}
      \caption{General Registers.}
      \label{general_operation_csrs_tab:is}
    \end{table}
\clearpage
\fi

\ifdefined\SECTIONCLEARPAGE
\clearpage
\fi
\section{Usage}

\subsection{Instantiation}
\label{sec:inst}
% SPDX-FileCopyrightText: 2024 IObundle
%
% SPDX-License-Identifier: MIT

Figure~\ref{fig:inst} illustrates how to instantiate the IP core and, if applicable, the
required external blocks.

\begin{figure}[!htbp]
    \centerline{\includegraphics[width=14cm]{inst.pdf}}
    \vspace{0cm}\caption{Core Instance and Required Surrounding Blocks}
    \label{fig:inst}
\end{figure}

% SPDX-FileCopyrightText: 2024 IObundle
%
% SPDX-License-Identifier: MIT

The figure shows a three-level memory hierarchy comprising L1 and L2 caches and
a memory module. The Host drives the L1 cache through its front-end NPI
interface (the user is free to develop other interfaces).  The L1 and L2 caches
are connected using another NPI interface since IOb-Cache's front-end interface
only supports this.

The {\tt wtb\_empty\_in} and {\tt wtb\_empty\_out} signals form a chain from the
L1's front-end to the L2's back-end. As explained in the description of these
signals, this chain ensures that the user's core knows that all write-through
buffers across the cache hierarchy are empty. Note that the L1's {\tt
  wtb\_empty\_out} signal is floating because the Host uses the cache controller
to query the write-through buffer status. The L2's {\tt wtb\_empty\_in} is tied
high as L2 is the last cache in the hierarchy, and there are no more
write-through buffers to its right-hand side.

The {\tt invalidate\_in} and {\tt invalidate\_out} signals form another chain that
ensures that the data in the whole cache hierarchy is invalidated, as explained
in these signal's descriptions. Note that the L1's {\tt invalidate\_in} signal is
tight to low as L1 is invalidated via the cache controller by writing to the
respective address. The L2's {\tt invalidate\_out} signal is floating because L2
is the last cache in the hierarchy, and there are no more caches to invalidate.

Finally, L2 is connected to a memory module, and one can choose between NPI or
AXI4 interfaces. In practice, most memory modules have a standard interface such
as AXI4, which is themost common choice, although one may choose NPI in less
usual simulation or FPGA prototyping scenarios.





\subsection{Simulation}
\label{sec:tbbd}
% SPDX-FileCopyrightText: 2024 IObundle
%
% SPDX-License-Identifier: MIT

The provided testbench uses the core instance described in
Section~\ref{sec:inst}. A high-level block diagram of the testbench is shown in
Figure~\ref{fig:tbbd}. The testbench is organized in a modular fashion, with
each test described in a separate file. The test suite consists of all the test
case files to make adding, modifying, or removing tests easy.

\begin{figure}[!htbp]
    \centerline{\includegraphics[width=14cm]{tbbd.pdf}}
    \vspace{0cm}\caption{Testbench Block Diagram}
    \label{fig:tbbd}
\end{figure}

% SPDX-FileCopyrightText: 2024 IObundle
%
% SPDX-License-Identifier: MIT

The above paragraph describes a desirable simulation setup, but IOb-Cache's
simulation environment still lacks a modular simulation structure. Currently,
only a set of primary non-pipelined write followed by read tests is
implemented. However, IOb-Cache has been thoroughly verified in-system, with two
cache levels. Various open-source RISC-V processors have proven that IOb-Cache
works well: PicoRV32, SSRV, VexRISCV, and DarkRV.
\clearpage


\subsubsection*{System-level Simulation}

Upon request, simulation files to run the core embedded in a RISC-V system can
be provided. The core is exercised in various modes by the RISC-V processor,
using a bare-metal software program written in the C programming language.


\ifdefined\ASICSYNTH
\subsection{ASIC Synthesis}
\label{sec:synth}
% SPDX-FileCopyrightText: 2024 IObundle
%
% SPDX-License-Identifier: MIT

A {\tt .tcl} script is provided for the Cadence Genus synthesis tool. The script
reads the technology files, compiles and elaborates the design, and proceeds to
synthesize it. Timing constraints file(s) are provided.

After synthesis, reports on silicon area usage, power consumption, and timing
closure are generated. A post-synthesis Verilog file is created, which can be
used in post-synthesis simulation. The IP core has been synthesized without the
FIFO memories, which are treated as black boxes. RAM blocks are not included as
they are technonoly dependent. The user is free to generate and add the required
memories, or request this service from the IP core provider.

\fi

\ifdefined\FPGACOMP
\subsection{FPGA Compilation}
\label{sec:fpga}
% SPDX-FileCopyrightText: 2024 IObundle
%
% SPDX-License-Identifier: MIT

The Quartus and Vivado FPGA compilation toolchains are supported via {\tt .tcl}
scripts invoked by a Makefile. The script compiles and elaborates the design for
a given set of target FPGA boards. Additional boards can be supported by
following the flow provided.

The core is instantiated within a Verilog wrapper, which depends on the
specifics of the FPGA device and the FPGA board. Timing constraints file(s) are
provided.

After compilation, reports on FPGA resource usage, power consumption, and timing
closure are generated. A post-synthesis Verilog file is created, which can be
used in post-synthesis simulation.

% SPDX-FileCopyrightText: 2024 IObundle
%
% SPDX-License-Identifier: MIT

Add here any specific ip core description.


\subsubsection*{System-level FPGA Run}

Upon request, files to run the core on the supported FPGA boards can be
provided. The core is embedded in a RISC-V system and exercised in various
modes, using a bare-metal software program written in the C programming
language.

\fi

\ifdefined\RESULTS
\ifdefined\SECTIONCLEARPAGE
\clearpage
\fi
\section{Implementation Results}
\label{sec:results}
% SPDX-FileCopyrightText: 2024 IObundle
%
% SPDX-License-Identifier: MIT

\ifdefined\FPGA
\subsection{FPGA}
% SPDX-FileCopyrightText: 2024 IObundle
%
% SPDX-License-Identifier: MIT

\ifdefined\RESULTS
\ifdefined\FPGA

\ifdefined\AMD
\begin{table}[H]
\centering
\begin{tabular}{|l|r|}
\hline
\rowcolor{iob-green}
\textbf{Resource}  & \textbf{Used} \\
\hline
\hline
\input vivado.tex 
\end{tabular}
\caption{AMD Kintex Ultrascale FPGAs}
\label{tab:fpga_amd}
\end{table}
\fi

\ifdefined\INTEL
\begin{table}[H]
\centering
\begin{tabular}{|l|r|}
\hline
\rowcolor{iob-green}
\textbf{Resource}  & \textbf{Used} \\
\hline
\hline
\input quartus.tex 
\end{tabular}
\caption{Intel Cyclone V GT FPGAs}
\label{tab:fpga_intel}
\end{table}
\fi

\fi
\fi

\fi


\ifdefined\ASIC
\subsection{ASIC}
% SPDX-FileCopyrightText: 2024 IObundle
%
% SPDX-License-Identifier: MIT

\ifdefined\ASIC
\begin{table}[H]
\centering
\begin{tabular}{|c|c|c|c|}
\hline
\rowcolor{iob-green}
\textbf{Area} ($\SI{}{\micro\meter\squared}$)  & \textbf{Gates}  & \textbf{FFs}  & \textbf{Frequency} ($\SI{}{\mega\hertz}$)\\
\hline
\hline
\input asic.tex
\end{tabular}
\caption{ASIC results for node UMC 130nm.}
\label{tab:asic_results}
\end{table}
\fi

\fi

\fi

\ifdefined\CUSTOM
\input{custom}
\fi

\end{document}
